\documentclass[10pt,a4paper]{amsart}
\usepackage[margin=19mm]{geometry}
\usepackage{amsmath,amssymb,mathtools}
\usepackage[scr=rsfs,cal=euler]{mathalfa}
\usepackage{xfrac}
\usepackage[unicode,hidelinks,bookmarks=false]{hyperref}
\newcommand{\ie}{i.e.\ }
\newcommand{\cf}{cf.\ }
\newcommand{\resp}{respectively\ }
\newcommand{\Subsection}{\subsection}
\providecommand{\nobreakdash}{\nobreak\hbox{-}}
\setlength{\emergencystretch}{2em}
\multlinegap0pt
\usepackage{bm}
\usepackage{bbm}
\usepackage{dsfont}
\usepackage{xspace}
\usepackage{cancel}
\usepackage{mathtools}
\usepackage{amsthm}
\makeatletter
\@ifundefined{theorem}{\newtheorem{theorem}{Theorem}}{}
\@ifundefined{lemma}{\newtheorem{lemma}[theorem]{Lemma}}{}
\makeatother
\newcommand{\Om}{{\Omega}}
\newcommand{\pa}{\partial}
\newcommand{\qie}{{\quad \mbox{i.e.}, \quad}}
\newcommand{\qin}{{\quad \mbox{in} \quad}}
\newcommand{\qon}{{\quad \mbox{on} \quad}}
\newcommand{\ve}{{\varepsilon}}
\newcommand{\vp}{{\varphi}}
\title[On the Convergence of Solutions for the Ginzburg-Landau Equa]{On the Convergence of Solutions for the Ginzburg-Landau Equation and System}
\author{Rejeb Hadiji, Jongmin Han}
\date{Source: arXiv:2509.09231v1 (CC BY 4.0)}
\begin{document}
\maketitle

\noindent\emph{Source and licence.} On the Convergence of Solutions for the Ginzburg-Landau Equation and System. Version arXiv:2509.09231v1 (CC BY 4.0), distributed under the Creative Commons Attribution 4.0 International licence, as recorded in the arXiv licence metadata for this exact version and shown on the abstract page https://arxiv.org/abs/2509.09231v1. Full rights notice: licenses/arxiv-2509.09231-CC-BY-4.0.txt.\par\smallskip

\noindent\textit{Editorial scope.} Counted unit: the Lemma whose source label is \texttt{lem:limsup} in the article (source0.tex lines 879--885, source label \texttt{lem:limsup}), delivered together with the article's own complete original proof of it (source0.tex lines 887--947, one \texttt{proof} environment of 61 lines). No mathematics is written, repaired or re-derived: statement and proof are the article's own text line for line. Required context retained uncounted: eq:GL (lines 316--324); lem:unif conv (lines 822--828); eq:nabla u L2 (lines 825--827); eq:nabla form (lines 844--846); eq-rho-zeta-1 (lines 848--852); eq-rho-zeta-2 (lines 848--852). Every definition, display and label of those lines is retained unchanged. Omitted from this package and not redistributed: 1--315, 325--821, 828--843, 847--847, 853--878, 886--886, 948--1430. Nothing inside a retained excerpt is omitted or altered: every retained body is delivered exactly as the article's lines are written, so no omission receipt of any kind accompanies this package. Citations: the retained excerpts carry no citation command at all, so no bibliography entry of the article is transcribed and no reference is redistributed. Reference adaptation: none was needed and none was made; every reference inside the delivered unit resolves inside it, so no unresolved reference or citation is produced. Printed slips kept literal: no printed word, symbol, hypothesis or number of the counted statement or of its proof was altered. Layout and class: the article's own class is \texttt{amsart} with packages including amsmath, amsthm, amsfonts, amssymb, amsopn, color, lineno, xy, esint, enumerate, mathrsfs, lipsum, pgf, tikz; the wrapper is the lane's pinned \texttt{amsart} editorial wrapper at 10pt on A4, which restates the theorem-like environments the retained text needs and adds the article's own macro definitions that the retained text needs (\texttt{\textbackslash Om}, \texttt{\textbackslash pa}, \texttt{\textbackslash qie}, \texttt{\textbackslash qin}, \texttt{\textbackslash qon}, \texttt{\textbackslash ve}, \texttt{\textbackslash vp}), with extra wrapper packages (bm, bbm, dsfont, xspace, cancel, mathtools). The delivered numbering is the unit's own. Assets: no retained span contains a figure environment, a graphics-inclusion command, a PDF-page inclusion, a table or a TikZ picture; no image file is copied into this package, the package declares no asset dependency and no image file is redistributed with it. Licence: version arXiv:2509.09231v1 of this article is distributed under the Creative Commons Attribution 4.0 International licence, as recorded in the arXiv licence metadata for this exact version and shown on the abstract page https://arxiv.org/abs/2509.09231v1. A full rights notice is delivered at licenses/arxiv-2509.09231-CC-BY-4.0.txt. Characters: every retained excerpt is pure ASCII, so the delivered engine, XeLaTeX with its default Latin Modern text font, reports no missing character.
\begin{equation}
\label{eq:GL}
\left\{
\begin{aligned}
-\Delta u& = \frac{1}{\ve^2} u (1-|u|^2 ) \qin \Omega, \\
u&=g \qquad  \qquad \quad \qon \pa\Omega.
\end{aligned}
\right.
\end{equation}

\begin{lemma}\label{lem:unif conv}
Let  $u_\ve$ be a solution for  \eqref{eq:GL}$_\ve$.
If $|u_\ve| \to 1 $ uniformly on $ \overline\Om $, then
\begin{equation}\label{eq:nabla u L2}
 \int_\Om  |\nabla u_\ve|^2  dx \le 2  \int_\Om  |\nabla u_0|^2  dx .
  \end{equation}
\end{lemma}

\begin{equation}\label{eq:nabla form}
 |\nabla u_\varepsilon|^2 =    |\nabla \rho_\varepsilon|^2 +  \rho_\varepsilon^2 |\nabla \zeta_\varepsilon|^2
\end{equation}

\begin{align}\label{eq-rho-zeta-1}
\operatorname{div}\left( \rho_\varepsilon^2 \nabla \zeta_\varepsilon \right) &= 0 \qin \Omega, \\
\label{eq-rho-zeta-2}
-\Delta \rho_\ve + \rho_\ve\big| \nabla  \zeta_\varepsilon  \big|^2 & = \frac{1}{\ve^2}\rho_\ve(1- \rho_\ve^2) \qin \Om.
\end{align}

\begin{lemma}\label{lem:limsup}
Let  $u_\ve$ be a solution for  \eqref{eq:GL}$_\ve$.
If $|u_\ve| \to 1 $ uniformly on $\overline\Om $, then
\begin{equation}\label{eq:limsup}
 \limsup_{\ve \to \infty} \int_\Om |\nabla u_\ve|^2  dx \le  \int_\Om |\nabla u_0|^2  dx.
\end{equation}
\end{lemma}

\begin{proof}
Let us assume the contrary.
Then,  there exists a constant $C_2>0$ and a subsequence, still denoted by $u_\ve$, such that
\begin{equation}\label{eq:larger than B1}
\int_\Omega |\nabla \vp|^2 dx=\frac12 \int_\Om |\nabla u_0|^2 dx< C_2 \le  \frac12 \int_\Om |\nabla u_\ve|^2 dx.
\end{equation}
 Since  $|u_\varepsilon| \to 1$ uniformly on $\overline\Omega$, we may keep the notations in the proof of Lemma \ref{lem:unif conv}.
Given   $\delta \in (0,\frac14)$, if $\ve$ is small enough, then
\begin{equation}\label{eq:delta chosen}
\frac{1}{2} < \rho_\ve <\rho_\ve^2 +\delta  \qie \frac{1+\sqrt{1-4\delta} }{2} <\rho_\ve <1.
\end{equation}
Let
\[ \psi_\varepsilon = \zeta_\varepsilon - \varphi.
\]
Then, by \eqref{eq:nabla form} and \eqref{eq:larger than B1},
\begin{equation}\label{eq:nabla u formula}
\begin{aligned}
C_2 \le  \frac{1}{2} \int_\Omega |\nabla u_\varepsilon|^2  dx
&=    \frac{1}{2} \int_\Omega |\nabla \rho_\varepsilon|^2  dx
+\frac{1}{2} \int_\Omega \rho_\varepsilon^2 \big|\nabla (\vp +\psi_\ve) \big|^2 dx.
\end{aligned}
\end{equation}
We rewrite \eqref{eq-rho-zeta-1} and \eqref{eq-rho-zeta-2} as 
\begin{align}\label{eq-rho-psi-1}
\operatorname{div}\left( \rho_\varepsilon^2 \nabla( \varphi + \psi_\varepsilon ) \right) &= 0 \qin \Omega, \\
\label{eq-rho-psi-2}
-\Delta \rho_\ve + \rho_\ve\big| \nabla (\varphi + \psi_\varepsilon )\big|^2 & = \frac{1}{\ve^2}\rho_\ve(1- \rho_\ve^2) \qin \Om.
\end{align}
Multiplying \eqref{eq-rho-psi-2} by $ \rho_\varepsilon - 1 $  and  integrating it over $ \Omega $, and using the boundary
condition $ \rho_\varepsilon  =1 $ on $\partial\Omega$, we obtain
\begin{equation}\label{eq-rho-psi-2-1}
\begin{aligned}
&\frac12\int_\Omega |\nabla \rho_\varepsilon|^2  dx +
\frac12 \int_\Omega\rho_\ve^2 | \nabla (\varphi + \psi_\varepsilon )|^2 dx - \frac12 \int_\Omega \rho_\ve | \nabla (\varphi + \psi_\varepsilon )|^2 dx \\
=~& \frac{1}{2\ve^2}\int_\Omega \rho_\ve( \rho_\varepsilon - 1)(1- \rho_\ve^2) dx  ~\leq~ 0.
\end{aligned}
\end{equation}
Then,  from   \eqref{eq:delta chosen},  \eqref{eq:nabla u formula} and \eqref{eq-rho-psi-2-1}, it follows that
\begin{equation}\label{eq-C2-final}
\begin{aligned}
C_2 ~&\le ~   \frac12
\int_\Omega\rho_\ve \big| \nabla (\varphi + \psi_\varepsilon )\big|^2 dx
 ~ \leq~   \frac12   \int_\Omega (\rho_\varepsilon^2+\delta)  \big| \nabla (\varphi + \psi_\varepsilon )\big|^2  dx \\
&  \le ~   \frac12   \int_\Omega  \rho_\varepsilon^2   \big( | \nabla  \varphi |^2 +2 \nabla  \varphi \cdot  \psi_\varepsilon +|\nabla \psi_\varepsilon|^2  \big)\,  dx + \frac12 \delta    \| \nabla \zeta_\ve \|_2^2 \\
&  \le ~  \frac12   \int_\Omega     | \nabla  \varphi |^2 +\frac12   \int_\Omega  \rho_\varepsilon^2   \big(2 \nabla  \varphi \cdot  \psi_\varepsilon +|\nabla \psi_\varepsilon|^2  \big)\,  dx  +   \frac12 \delta    \| \nabla \zeta_\ve \|_2^2 .
\end{aligned}
\end{equation}
Multiplying \eqref{eq-rho-psi-1} by \( \psi_\varepsilon \), integrating it over \( \Omega \), and using the boundary condition \( \psi_\varepsilon = 0 \) on \( \partial\Omega \), we obtain
\begin{align}\label{eq-rho-psi-1-1}
  \int_\Omega \rho_\varepsilon^2 \left| \nabla \psi_\varepsilon \right|^2 \, dx  +\int_\Omega \rho_\varepsilon^2 \nabla \varphi \cdot \nabla \psi_\varepsilon \, dx =0.
\end{align}
Furthermore, by \eqref{eq:nabla u L2} and \eqref{eq:delta chosen},
\begin{equation}\label{eq:nabla zeta}
   \| \nabla \zeta_\ve \|_2^2 \le 4 \int_\Om \rho_\ve^2 | \nabla \zeta_\ve|^2 dx \le 8 \int_\Om |\nabla u_0|^2 dx.
\end{equation}
Hence, by   \eqref{eq-C2-final},  \eqref{eq-rho-psi-1-1} and \eqref{eq:nabla zeta}, we are led to
\begin{align*}
 0< C_2 -   \frac12   \int_\Omega     | \nabla  \varphi |^2 \le   -   \int_\Omega \rho_\varepsilon^2 \left| \nabla \psi_\varepsilon \right|^2 \, dx+  4 \delta  \|\nabla u_0 \|_2^2 \le   4 \delta   \|\nabla u_0|\|_2^2.
\end{align*}
Letting $\delta \to 0$, we arrive at  a contradiction.
\end{proof}

\end{document}
